\section*{Hoofdstuk \ref{ch:b2:pt-paths} --- Thermobarometrie en P--T--t-paden}

\begin{solution}{exo:b2:pt-paths:1}
(a) A chlorietzone, B biotietzone, C granaatzone, D staurolietzone, E kyanietzone,
F sillimanietzone. (b) Vijf: de isograden van biotiet, granaat, stauroliet, kyaniet en sillimaniet,
elk tussen twee opeenvolgende monsters, met de graad stijgend van A naar F. (c) Een isograad
zegt alleen dat een reactie is verlopen; een veldgradiënt vraagt drukken en temperaturen, die
komen uit gekalibreerde reacties of uit thermobarometrie van de piekassociatie van elk
gesteente.
\end{solution}

\begin{solution}{exo:b2:pt-paths:2}
(a) $(X_{\text{Mg}}/X_{\text{Fe}})_{\text{grt}} = 0.10/0.80 = 0.125$,
$(X_{\text{Mg}}/X_{\text{Fe}})_{\text{bt}} = 0.45/0.50 = 0.900$; $K_D = 0.1389$, $\ln K_D =
-1.9741$. (b) $3R = 24.942$; noemer $20.19 + 24.942\times1.9741 = 69.428$; bij
\qty{5}{kbar} $T = 54185/69.428 = \qty{780.5}{K} = \qty{507}{\celsius}$; bij
\qty{7}{kbar} $T = (52630 + 2177)/69.428 = \qty{789.4}{K} = \qty{516}{\celsius}$.
(c) Twee kilobar veranderen haar maar met \qty{9}{K}: een thermometer.
\end{solution}

\begin{solution}{exo:b2:pt-paths:3}
(a) $z = P/(\bar\rho g) = 6.0\times10^8/(2800\times9.81) = \qty{21.8}{km}$.
(b) $(580 - 10)/21.8 = \qty{26}{\celsius\per\km}$. (c) Bij \qty{580}{\celsius} ligt de lijn
op $3.67 + 0.8\times2.10 = \qty{5.35}{kbar}$; bij \qty{6.0}{kbar} ligt het gesteente aan de
hogedrukkant: kyaniet.
\end{solution}

\begin{solution}{exo:b2:pt-paths:4}
(a) Reactanten $2\times140.693 + 3\times18.015 = \qty{335.431}{g}$; producten $277.112 +
58.320 = \qty{335.432}{g}$: in evenwicht op de afronding na. (b) $2\times43.65 =
\qty{87.30}{cm^3}$ wordt $107.50 + 24.63 = \qty{132.13}{cm^3}$: $+\qty{51.4}{\percent}$.
(c) Olivijn $281.386/87.30 = \qty{3.223}{g.cm^{-3}}$; producten $335.432/132.13 =
\qty{2.539}{g.cm^{-3}}$. (d) Het gesteente zwelt met de helft van zijn volume op en wordt lichter; het
slaat ook water op in zijn mineralen.
\end{solution}

\begin{solution}{exo:b2:pt-paths:5}
(a) Gemiddelden $-1.206$ bij \qty{1072.15}{K} en $-1.782$ bij \qty{823.15}{K}. Met $1/T =
9.3271\times10^{-4}$ en $1.21485\times10^{-3}$ is $a = (-1.206 + 1.782)/(9.3271\times10^{-4}
- 1.21485\times10^{-3}) = -2041.5$ K en $b = -1.206 + 2041.5/1072.15 = 0.6982$.
(b) $\Delta_r H + 2070\times0.311 = 2041.5\times24.942 = \qty{50920}{J}$, dus $\Delta_r H =
\qty{50276}{J}$; $\Delta_r S = 0.6982\times24.942 = \qty{17.41}{J.K^{-1}}$. (c) Bij
\qty{5}{kbar}: voor $K_D = 0.18$ 588.1 tegen \qty{587.5}{\celsius}; voor $0.10$ 419
tegen \qty{425}{\celsius}; voor $0.30$ 819 tegen \qty{806}{\celsius}. Binnen het bereik
van de experimenten komen de twee kalibraties overeen, hoewel hun constanten verschillen met
\qty{2.4}{kJ} en \qty{2.8}{J.K^{-1}}; de fouten van $\Delta_r H$ en $\Delta_r S$
compenseren elkaar daar en groeien daarbuiten. Een kalibratie mag niet ver
buiten de omstandigheden van haar experimenten worden gebruikt.
\end{solution}

\begin{solution}{exo:b2:pt-paths:6}
(a) Uitwisseling: $(113.12 + 154.30) - (115.32 + 149.65) = \qty{+2.45}{cm^3} =
\qty{0.245}{J.bar^{-1}}$, tegen 0.311. GASP: $125.28 + 2\times44.15 + 22.69 -
3\times100.79 = \qty{-66.10}{cm^3} = \qty{-6.61}{J.bar^{-1}}$, tegen $-6.21$, de waarde
van de kalibratie bij \qty{800}{\celsius} en \qty{10}{kbar}, waar thermische uitzetting en
samendrukking de volumes veranderen. In beide gevallen is de GASP-verandering twintig tot dertig keer
die van de uitwisseling. (b) $\partial T/\partial P = 0.311/62.96 = \qty{4.9e-3}{K/bar}$:
\qty{4.9}{K} per kilobar. (c) De helling van de lijn is $22.80 + 8.314\times3\ln 0.20/6.21
= 22.80 - 6.46 = \qty{16.3}{bar/K}$; \qty{17}{K} verschuift de druk met
\qty{280}{bar}.
\end{solution}

\begin{solution}{exo:b2:pt-paths:7}
(a) $P_0 = 22.80\times600 - 1093 = \qty{12587}{bar}$; $3\ln(0.08/0.30) = -3.9653$; $RT =
8.314\times873.15 = \qty{7259.4}{J}$; $P = 12587 + 7259.4\times(-3.9653)/6.21 = 12587 -
4635 = \qty{7952}{bar}$. (b) $3\ln(0.08/0.35) = -4.4277$, $P = 12587 - 5176 =
\qty{7411}{bar}$: \qty{0.54}{kbar} minder, dus het plagioklaas moet zorgvuldig worden geanalyseerd.
(c) $7952\times10^5/(2800\times9.81) = \qty{29.0}{km}$.
\end{solution}

\begin{solution}{exo:b2:pt-paths:8}
(a) $z = P/(\bar\rho g)$: \qty{118.9}{km}, \qty{34.0}{km} en \qty{15.3}{km}.
(b) $(118.9 - 34.0)/2.2 = \qty{38.6}{km/Myr} = \qty{3.9}{cm.yr^{-1}}$ en $(34.0 -
15.3)/1.1 = \qty{17.0}{km/Myr} = \qty{1.7}{cm.yr^{-1}}$, iets boven de gepubliceerde
3.4 en 1.6: hun omrekening van druk naar diepte geeft iets kleinere
diepteverschillen. (c) $\sigma = \sqrt{0.9^2 + 0.9^2} = \qty{1.3}{Myr}$, dus $2.2 \pm 1.3$
miljoen jaar en een snelheid tussen $84.9/3.47 = 24$ en $84.9/0.93 = \qty{91}{km/Myr}$,
2.4 tot \qty{9.1}{cm.yr^{-1}}: hoe dan ook snel. (d) Zelfs de laagste snelheid is ongeveer vijftig keer
de \qty{0.5}{mm.yr^{-1}} van het model: erosie alleen kan dit niet.
\end{solution}

\begin{solution}{exo:b2:pt-paths:9}
(a) $\partial T/\partial\ln K_D = 3RT/(\Delta_r S - 3R\ln K_D) = 24.942\times860.6/62.96 =
\qty{341}{K}$; $\partial T/\partial P = 0.311/62.96 = \qty{4.9}{K}$ per kilobar.
(b) $\sigma_T = \sqrt{(341\times0.05)^2 + 4.9^2} = \sqrt{17.0^2 + 4.9^2} =
\qty{17.7}{K}$. (c) De analysefout op $K_D$ overheerst; die alleen is twee derde van
de \qty{25}{\celsius} van de beste kalibraties, die ook de fouten van de
kalibratie zelf omvatten.
\end{solution}

\begin{solution}{exo:b2:pt-paths:10}
(a) Granaat $0.10/0.75 = 0.1333$, biotiet $0.42/0.52 = 0.8077$, $K_D = 0.1650$, $\ln K_D =
-1.8018$, noemer $20.19 + 24.942\times1.8018 = 65.130$; $3\ln(0.07/0.25) = -3.8189$.
Bij \qty{6000}{bar}: $T = 54496/65.130 = \qty{836.7}{K} = \qty{563.6}{\celsius}$, daarna $P
= 22.80\times563.6 - 1093 + 8.314\times836.7\times(-3.8189)/6.21 = \qty{7479}{bar}$.
Bij \qty{7479}{bar}: $T = \qty{570.6}{\celsius}$, $P = \qty{7603}{bar}$. Bij
\qty{7603}{bar}: $T = \qty{571.2}{\celsius}$, $P = \qty{7614}{bar}$; nog eens: $T =
\qty{571.3}{\celsius}$. Het gesteente kwam in evenwicht bij ongeveer \qty{571}{\celsius} en
\qty{7.6}{kbar}. (b) Bij \qty{571}{\celsius} ligt de lijn kyaniet $=$ sillimaniet op $3.67
+ 0.71\times2.10 = \qty{5.16}{kbar}$: kyaniet. (c) De barometer werd gekalibreerd voor
kyaniet; met sillimaniet verschillen de reactie en haar lijn, omdat de twee polymorfen niet
dezelfde gibbsenergie hebben.
\end{solution}

\begin{solution}{exo:b2:pt-paths:11}
(a) $\tau = L^2/(\pi^2\kappa)$: voor \qty{35}{km} $1.225\times10^9/(9.870\times10^{-6}) =
\qty{1.24e14}{s} = \qty{3.9}{Myr}$; voor \qty{70}{km} vier keer meer,
\qty{15.7}{Myr}. (b) Zeventig miljoen jaar is vier tot achttien keer deze tijden: warmte
verplaatst zich veel sneller door de verdikte korst dan erosie haar wegneemt. (c) Het gesteente
bereikt zijn grootste diepte meteen, bij het stapelen; de warmte komt in de daaropvolgende
tientallen miljoenen jaren aan, terwijl de erosie het langzaam omhoogbrengt: de temperatuur blijft
een tijdlang stijgen tijdens de decompressie. (d) Erosie in zeven miljoen jaar zou
het gesteente omhoogbrengen voordat geleiding het had verhit: een lagere thermische piek, vroeg
bereikt, bij hogere druk, en een koud pad terug naar het oppervlak.
\end{solution}

\begin{solution}{exo:b2:pt-paths:12}
(a) De korst weegt $35000\times2800\times9.81 = \qty{0.961}{GPa}$; de resterende
\qty{1.839}{GPa} vragen $1.839\times10^9/(3300\times9.81) = \qty{56.8}{km}$ mantel: diepte
\qty{91.8}{km}, zoals de \qty{90}{km} of meer van de oorspronkelijke studie. (b) $(4.0 -
0.961)\times10^9/(3300\times9.81) = \qty{93.9}{km}$, diepte \qty{128.9}{km}.
(c) $700/91.8 = \qty{7.6}{\celsius\per\km}$: een zeer hoge verhouding van druk tot
temperatuur, de reeks van subductiezones. (d) $(91.8 - 35)/34 = \qty{1.7}{Myr}$.
\end{solution}

\begin{solution}{pb:b2:pt-paths:1}
\noindent\textbf{Deel I.}

\textbf{1.} Granaatkern $0.080/0.760 = 0.1053$, biotiet $0.40/0.58 = 0.6897$; $K_D =
0.1527$, $\ln K_D = -1.8793$.

\textbf{2.} Rand $0.120/0.720 = 0.1667$, matrixbiotiet $0.46/0.53 = 0.8679$; $K_D =
0.1921$, $\ln K_D = -1.6497$.

\textbf{3.} $3R = 24.942$; teller $52630 + 0.311\times6000 = 54496$; noemer $20.19
+ 24.942\times1.8793 = 67.064$; $T = \qty{812.6}{K} = \qty{539.5}{\celsius}$.

\textbf{4.} Noemer $20.19 + 24.942\times1.6497 = 61.337$; $T = 54496/61.337 =
\qty{888.5}{K} = \qty{615.3}{\celsius}$.

\textbf{5.} $0.311\times1000/67.064 = \qty{4.6}{K}$ voor de kern en $0.311\times1000/61.337
= \qty{5.1}{K}$ voor de rand.

\textbf{6.} Een thermometer geldt voor mineralen die samen in evenwicht waren. Het
biotiet dat in de kern is ingesloten, werd van de rest van het gesteente afgesloten toen de kern groeide;
het matrixbiotiet wisselde het laatst uit met de rand. De kern koppelen aan de matrix zou
twee momenten door elkaar halen.

\noindent\textbf{Deel II.}

\textbf{7.} Kern $22.80\times539.5 - 1093 = \qty{11208}{bar}$; rand $22.80\times615.3 -
1093 = \qty{12936}{bar}$.

\textbf{8.} $3\ln(0.050/0.25) = -4.8283$; $RT = 8.314\times812.65 = \qty{6756.4}{J}$; $P =
11208 - 6756.4\times4.8283/6.21 = 11208 - 5253 = \qty{5955}{bar}$.

\textbf{9.} $3\ln(0.060/0.28) = -4.6213$; $RT = 8.314\times888.45 = \qty{7386.6}{J}$; $P
= 12936 - 5497 = \qty{7439}{bar}$.

\textbf{10.} Bij \qty{7439}{bar}: $T = (52630 + 2314)/61.337 = \qty{895.8}{K} =
\qty{622.6}{\celsius}$; bij die temperatuur $P = 13102 - 5542 = \qty{7560}{bar}$; daarna $T
= \qty{623.2}{\celsius}$. De rand groeide bij ongeveer \qty{623}{\celsius} en \qty{7.56}{kbar}.

\textbf{11.} Bij \qty{5954}{bar} geldt $T = \qty{539.2}{\celsius}$ en $P = \qty{5950}{bar}$: de
kern groeide bij ongeveer \qty{539}{\celsius} en \qty{5.95}{kbar}. De eerste schatting,
\qty{6}{kbar}, zat er al dicht bij, en de thermometer verschuift maar \qty{4.6}{K} per kilobar.

\textbf{12.} Kern $5950\times10^5/(2800\times9.81) = \qty{21.7}{km}$; rand
$7560\times10^5/27468 = \qty{27.5}{km}$.

\textbf{13.} Bij \qty{539}{\celsius} ligt de lijn op $3.67 + 0.39\times2.10 =
\qty{4.49}{kbar}$, onder \qty{5.95}{kbar}: kyaniet, consistent met het kyaniet dat in de
kern is ingesloten. Bij \qty{623}{\celsius} ligt ze op $5.77 + 0.23\times2.09 = \qty{6.25}{kbar}$,
onder \qty{7.56}{kbar}: weer kyaniet.

\noindent\textbf{Deel III.}

\textbf{14.} $+\qty{84}{\celsius}$ en $+\qty{1.6}{kbar}$ (ongeveer \qty{6}{km} dieper): het
gesteente werd begraven en verhit, op zijn \omterm{def:b2:pt-paths:paths}{prograde pad}.

\textbf{15.} Bij \qty{623}{\celsius} of meer, onder \qty{6.25}{kbar}: nadat de rand was gegroeid,
verloor het gesteente minstens \qty{1.3}{kbar}, ongeveer \qty{5}{km} bovenliggend gesteente, terwijl het minstens
even heet bleef. Het werd geëxhumeerd terwijl het nog heet was.

\textbf{16.} Chloriet is een waterhoudend mineraal van lage graad: het legt het late \omterm{def:b2:pt-paths:paths}{retrograde
pad} vast, afkoeling op geringe diepte, met water.

\textbf{17.} Begraving en verhitting tot de rand, decompressie zonder afkoeling, daarna afkoeling:
getekend met de temperatuur naar rechts en de druk omhoog loopt het pad met de klok mee. De
drukpiek kwam eerst, bij de rand; de thermische piek tegelijk of later, tijdens de
decompressie.

\textbf{18.} Ja: ook het model bereikt eerst zijn drukpiek en zijn thermische piek,
\qty{593}{\celsius} bij \qty{7.8}{kbar}, tijdens het opstijgen, dicht bij de omstandigheden
van de rand van de schist. Verdikking gevolgd door trage erosie verklaart het pad van de schist.

\noindent\textbf{Deel IV.}

\textbf{19.} $(623 - 405)/(412.0 - 398.0) = 218/14.0 = \qty{15.6}{\celsius}$ per miljoen
jaar.

\textbf{20.} Het gesteente koelde ongeveer anderhalf keer sneller af dan aangenomen. Een mineraal
dat sneller afkoelt, heeft minder tijd om zijn argon te verliezen, dus sluit het bij een hogere temperatuur:
de werkelijke sluitingstemperatuur ligt wat boven \qty{405}{\celsius}, en de sluiting
wat dieper.

\textbf{21.} $(405 - 10)/30 = \qty{13.2}{km}$.

\textbf{22.} $27.5 - 13.2 = \qty{14.3}{km}$ (met de onafgeronde diepten, 27.52 en 13.17,
\qty{14.35}{km}).

\textbf{23.} $\sqrt{1.5^2 + 1.0^2} = \qty{1.8}{Myr}$: $14.0 \pm \qty{1.8}{Myr}$, dus een snelheid
tussen $14.35/15.8 = 0.91$ en $14.35/12.2 = \qty{1.18}{km/Myr}$.

\textbf{24.} $14.35/14.0 = \qty{1.03}{km}$ per miljoen jaar, dat wil zeggen
\textbf{\qty[text-series-to-math=true]{1.0}{mm.yr^{-1}}}, ongeveer dertig keer trager dan
het eerste stadium van \omterm{def:b2:pt-paths:exhumation}{exhumatie} van de diepste gesteenten van de Alpen, en ongeveer tweemaal de erosiesnelheid
van het verdikkingsmodel.
\end{solution}
